\documentclass[10pt,a4paper]{article}

\usepackage[T1]{fontenc}
\usepackage[utf8]{inputenc}
\usepackage{lmodern}
\usepackage{geometry}
\usepackage{hyperref}
\usepackage{microtype}
\usepackage{graphicx}
\usepackage{amsmath}
\usepackage{tabularx}
\usepackage{array}
\usepackage{longtable}

\geometry{margin=2.5cm}
\setlength{\emergencystretch}{2em}
\newcolumntype{L}[1]{>{\raggedright\arraybackslash}p{#1}}
\newcolumntype{Y}{>{\raggedright\arraybackslash}X}
\hypersetup{
  colorlinks=true,
  linkcolor=black,
  citecolor=black,
  urlcolor=blue,
  pdftitle={ImpactLLM Methodology: Transparent Screening for LLM Inference and Training Impacts},
  pdfauthor={Arnault Pachot and Thierry Petit}
}


\title{ImpactLLM Methodology: Transparent Screening for LLM Inference and Training Impacts}
\author{Arnault Pachot \and Thierry Petit}
\date{March 14, 2026}

\begin{document}

\maketitle

\begin{abstract}
This paper describes the implementation of the ImpactLLM estimator. While the observatory publishes the reference table for 41 models, this companion manuscript focuses on the algorithms that convert natural-language application descriptions into inference and training carbon estimates.
\end{abstract}

\textbf{Keywords:} LLMs, source-linked evidence, screening method, natural language interface.
\section{Introduction}
This paper accompanies the \emph{ImpactLLM Observatory} and documents the proxies that produce the inference and training carbon numbers published there. The observatory now covers 41 models, and this companion paper explains how the interface extracts scenarios from natural-language descriptions and propagates them through the same bounded multi-factor screening proxies explained below. The goal is to keep those proxies auditable while keeping the application workflow fast and interpretable.

The remainder of this paper focuses on the technical design of the inference estimator, the training proxy, and the limits of those assumptions. The actual environmental benchmarks remain in the observatory, which references these methods for transparency.


\section{Inference Screening Method}
The main practical question addressed by the tool is straightforward: given a software feature using an LLM, can we produce an inference estimate that is useful for screening and comparison when direct provider telemetry is absent?

The current answer is deliberately limited. The estimator is \emph{inference-only}. It excludes model training, embodied impacts, application-side software consumption, and ancillary infrastructure from the displayed result. This is not because those dimensions are unimportant, but because the available inference anchors are structured enough to support a transparent screening method, whereas mixing them with broader lifecycle terms would blur the meaning of the result.

The design objective is also practical speed. Instead of asking users to fill a large technical form, the web interface starts from a natural-language description such as ``We use GPT-4o-mini for customer support, around 4{,}000 uses per month in France.'' A parser then maps this text to a compact scenario with model, request type, approximate tokens, usage volume, and country assumptions. The output is therefore fast to obtain, but every inferred parameter remains visible to the user and can be inspected or challenged.

The current market-model release follows five rules.
\begin{itemize}
\item It starts from an observed literature anchor rather than from provider claims. The retained prompt-energy anchor is the Gemini Apps median prompt energy reported by Elsworth et al.: $0.24$ Wh/prompt \cite{elsworth2025}.
\item It makes token volume explicit. Prompt compute is approximated from input and output tokens, with a larger weight for output generation than for input processing.
\item It does not scale on raw parameters alone. It constructs \emph{effective active parameters} by adjusting the target profile with context-window class, serving mode, modality support, and architecture notes.
\item It reports a bounded low-central-high interval rather than a single falsely precise point.
\item It derives carbon from retained energy and the electricity mix of the retained country, which is a contextual assumption rather than a model measurement.
\end{itemize}

Formally, let $E_a = 0.24$ Wh be the anchor value and $P_a = 180$ the active-parameter proxy retained for the anchor model. For a target model $t$, let $P_t$ denote active parameters, $T_{in}$ input tokens, and $T_{out}$ output tokens. We define weighted prompt-compute volume:
\begin{equation}
V_t = T_{in} + \omega T_{out},
\end{equation}
with $\omega = 1.8$ in the current release. The project reference scenario is:
\begin{equation}
V_{ref} = 1000 + 1.8 \times 550 = 1990.
\end{equation}

Effective active parameters are then defined by:
\begin{equation}
P^{eff}_{t,s} = P_t \times F^{ctx}_{t,s} \times F^{srv}_{t,s} \times F^{mod}_{t,s} \times F^{arch}_{t,s},
\end{equation}
where $s \in \{\mathrm{low}, \mathrm{central}, \mathrm{high}\}$ indexes the screening scenario. The final per-request energy estimate is:
\begin{equation}
\hat{E}_{t,s} = E_a \times \left(\frac{P^{eff}_{t,s}}{P_a}\right)^{\alpha_s}
\times \left(\frac{V_t}{V_{ref}}\right)^{\beta_s}.
\end{equation}
The current implementation uses $\alpha_{low}=0.85$, $\alpha_{central}=0.95$, $\alpha_{high}=1.05$, and $\beta_{low}=0.85$, $\beta_{central}=0.92$, $\beta_{high}=1.0$. Carbon is then derived from the retained country mix:
\begin{equation}
\hat{C}_{t,s,c} = \frac{\hat{E}_{t,s}}{1000} \times CI_c,
\end{equation}
where $CI_c$ is the carbon intensity of country $c$ in gCO$_2$e/kWh.

This method is not presented as a physical law of inference scaling. It is a traceable screening proxy anchored in observed prompt energy and explicit contextual assumptions. Its purpose is to provide a fast and inspectable estimate for comparative reasoning, not an audited declaration of real provider-side energy use.

\section{Training Screening Method}
The observatory also exposes training orders of magnitude for current market models. Here the uncertainty is even higher than for inference because published values are sparse, heterogeneous, and often reported only as aggregate emissions. A pure parameter-only scaling is therefore too brittle. The current release instead uses a second bounded proxy that combines retained parameter count with additional training priors.

The training proxy starts from literature anchors that directly report training CO$_2$e, and from training-energy reconstructions derived from those emissions when the source-country electricity mix is documented. For a target model $t$, the central training estimate is:
\begin{equation}
\hat{E}^{train}_{t,s} = E^{train}_a \times \left(\frac{P_t}{P_a}\right)^{\alpha_s}
\times \left(\frac{Tok_t}{Tok_a}\right)^{\beta_s}
\times F^{reg}_{t,s} \times F^{arch-tr}_{t,s} \times F^{hw}_{t,s},
\end{equation}
where $P_t$ is the retained parameter count, $Tok_t$ is a training-token prior, $F^{reg}$ is a training-regime prior, $F^{arch-tr}$ captures architecture and multimodality assumptions, and $F^{hw}$ is a hardware-class proxy. The current release uses a simple prior of 20 training tokens per retained parameter when no better public estimate is available, and defaults market models to a foundation-pretraining regime unless a narrower public indication exists.

This second proxy is not used in the application estimator shown to users, which remains inference-only. It is used in the observatory and the comparative tables to avoid suggesting that training impacts can be projected from parameter count alone. The result should still be read as a screening order of magnitude rather than as an audited declaration of model-development impact.

\section{Limitations}
The main limitations are methodological. First, the estimator depends on scarce literature anchors, so the uncertainty intervals rely on wide parameter- and token-exponent bounds to keep the screening range honest. Second, proprietary service configurations remain opaque, forcing the method to rely on contextual proxies for regime, architecture, and hardware. Third, training-level emissions are reconstructed from sparse reports and assumed token priors; more published training anchors would reduce that uncertainty. Finally, the natural-language parser must interpret ambiguous descriptions, which is why every extracted assumption is surfaced to the user along with the final estimate.

\section{Implementation Notes}
The prototype is implemented as a Python/Flask web application that builds on the observatory dataset stored in \\texttt{ImpactLLM/data/market_models.csv}. The front-end sends the user description to OpenAI for parsing, then queries the estimator to compute the inference and training proxies. The same dataset drives the observatory tables, which remain independent of the estimator logic.

\section{Conclusion}
The contribution of this manuscript is primarily methodological: it documents how ImpactLLM transforms natural-language descriptions into vetted inference and training carbon estimates. The observatory paper provides the reference table of 41 models, and this companion paper makes explicit how those numbers are derived and what assumptions underlie them.

\bibliographystyle{plain}
\bibliography{references_ImpactLLM}
\end{document}
